\section{Models and methods}
\label{sec:analysis}

\subsection{Point process models}

Let $(N_t)_{t\ge 0}$ denote the counting process associated with the
earthquake times $t_1 < \dots < t_N$ on $[0,T]$. All models are
specified through their conditional intensity
$\lambda(t \mid \mathcal{F}_t)$, where $\mathcal{F}_t$ is the history
up to time $t$.

The homogeneous Poisson model assumes a constant intensity
$\lambda(t) \equiv \lambda > 0$, so that the inter-event times are
i.i.d.\ $\mathrm{Exp}(\lambda)$. In contrast, Hawkes processes have the
form
\[
  \lambda(t) = \mu + \sum_{t_i < t} g(t - t_i, M_i),
\]
where $\mu > 0$ is a background rate and $g$ is a triggering kernel.
We consider four Hawkes specifications:
\begin{itemize}
  \item Exponential kernel (unmarked):
    $g(u,M) = \alpha e^{-\beta u} \mathbf{1}_{\{u>0\}}$.
  \item Power-law kernel (unmarked):
    $g(u,M) = K (c+u)^{-p} \mathbf{1}_{\{u>0\}}$.
  \item Marked exponential kernel:
    $g(u,M) = \alpha e^{\gamma (M-M_0)} e^{-\beta u} \mathbf{1}_{\{u>0\}}$.
  \item Marked power-law kernel:
    $g(u,M) = K e^{\gamma (M-M_0)} (c+u)^{-p} \mathbf{1}_{\{u>0\}}$.
\end{itemize}
All parameters are constrained to ensure positivity and stability
(e.g.\ $\alpha/\beta < 1$ in the exponential case).
All parameters are constrained to ensure positivity and stability
(e.g.\ $\alpha/\beta < 1$ in the exponential case, and $p \leq 5$
in the power-law specifications, as larger exponents yield a kernel
decay so rapid that the power-law and exponential specifications
become virtually indistinguishable).

\subsection{Maximum likelihood estimation}

For a given intensity model $\lambda_\theta(t)$ with parameter vector
$\theta$, the log-likelihood of the observed times is
\[
\ell(\theta)
  = \sum_{i=1}^N \log \lambda_\theta(t_i)
    - \int_0^T \lambda_\theta(t)\, \mathrm{d}t.
\]
For the homogeneous Poisson model the MLE is
$\hat{\lambda} = N / (T - t_0)$. For all Hawkes models we evaluate the
log-likelihood numerically and maximise it using the \texttt{scipy.minimize}
routine with a quasi-Newton method and simple box constraints on the
parameters.

\subsection{Residual analysis via time rescaling}

To assess absolute goodness of fit we use Ogata's time-rescaling
theorem. Given a fitted model with intensity $\lambda_{\hat\theta}$, we
compute the integrated intensity
$\Lambda_{\hat\theta}(t) = \int_{0}^{t} \lambda_{\hat\theta}(s)\,
\mathrm{d}s$ and define transformed times
$\tau_i = \Lambda_{\hat\theta}(t_i)$. Under correct specification, the
rescaled gaps $\Delta\tau_i = \tau_i - \tau_{i-1}$ should be i.i.d.\
$\mathrm{Exp}(1)$. We therefore apply Kolmogorov--Smirnov and
chi-squared tests against the $\mathrm{Exp}(1)$ law and compare the
empirical CDFs and densities of the $\Delta\tau_i$ to their theoretical
exponential counterparts. This allows us to check not only relative
fit, but also whether a given model can be regarded as an adequate
description of the data in absolute terms.
